\section{The index changes required at the floor 42}\label{ms:uniform-bridge}

The support proof in Section~\ref{ms:support} is uniform for $n\ge42$.
The other estimates also depend on the lower bound for $n$. We derive
the changes needed in concentration, confinement and normalization,
assuming the contact, phase-resource and local-comparison identities
stated in the preceding sections.

\subsection{Conditional moments and strict-gap concentration}
The concentration argument uses the antichain method of
Littlewood--Offord theory~\cite{Erdos1945} in its random-chain
form~\cite{Lubell1966}.

Let $n\ge N=42$ and use the same coordinates
$X=\ell(B_0+U-p)$ as in Section~\ref{ms:geometry-counts}. Suppose an
already proved support bound gives
\[
 \supp(B_0+U)\subset(-w_{\rm strip},1+w_{\rm strip}),\qquad
 U_0\le a_0,\quad U_1\ge-a_1,\quad
 a_0+a_1=d_Y\le d_0<w_{\rm strip}.
\]
The strip width is $23/40$ in the first left-panel contraction.
Any smaller value must be established before it is used in a subsequent
contraction. Assume $p\in[l,h]\subset(0,1/2)$.

Write $s=qw_0+pw_1=\lambda/n$, $w_i=\E U_i^2$. The exact identity
\eqref{ms:conditional-moments}, $|U_i|<w_{\rm strip}$ and $pq(w_0+w_1)\le qs$
give
\[
 \rho\le\rho_B(p)+c_\rho s
 \quad\hbox{if}\quad
 c_\rho>0,\quad c_\rho^2[l(1-l)]^3>[3(1-l)+w_{\rm strip}]^2.
\]
Indeed $|\E U_i^3|\le w_{\rm strip} w_i$,
$\ell^3\le[l(1-l)]^{-3/2}$, and the skeleton's third moment is at
most $\rho_B(p)$. The function $\rho_B(p)=(1-2p(1-p))/\sqrt{p(1-p)}$
decreases on $(0,1/2)$, so a closed noise band with ceiling
$\lambda_+$ gives
\[
 \rho\le\rho_B(l)+c_\rho\lambda_+/N.
\]
The first left contraction uses $c_\rho=30$. A later value such as $24$
is licensed only if its displayed gate holds with the actual strip.

Independently couple each pair of conditional residuals before choosing
its Bernoulli label. Changing a label from zero to one increases the
corresponding summand by at least $1-d_Y$. Conditional on the coupled
residuals, the inverse image of an interval of diameter strictly less
than $1-d_Y$ is therefore an antichain in the Boolean cube. For an
antichain $\mathcal A$ of subsets of $\{1,\ldots,j\}$, a uniformly
random permutation produces a chain of initial segments that intersects
$\mathcal A$ at most once. Taking expectations gives
\[
 \sum_{A\in\mathcal A}\binom{j}{|A|}^{-1}\le1.
\]
Under independent labels of parameter $p$,
\[
 \Pp(\mathcal A)=\sum_{A\in\mathcal A}
 \binom{j}{|A|}^{-1}\binom{j}{|A|}p^{|A|}q^{j-|A|}
 \le\max_k\binom jk p^kq^{j-k}.
\]
Conditioning back proves the strict-gap concentration bound for the
original sum, including any independent residual shift. The strict
width condition is necessary; a closed interval with width equal to the
increment can contain comparable vertices.

At floor $42$, the two deleted indices are at least $40$. Put
\[
 r_L=l-(1-l)/40,\qquad r_R=h+h/40,
 \qquad v_*=\min\{r_L(1-r_L),r_R(1-r_R)\}.
\]
The mode condition $|(k/j)-p|\le\max(p,q)/j$ puts every binomial mode
for $j\ge40$ in $[r_L,r_R]$. Require $0<r_L<r_R<1$.
For an interior mode with $r=k/j$, finite Stirling and nonnegativity of
the relative entropy give
\[
 \sqrt j\binom jk p^kq^{j-k}
 <\frac{\phi(0)}{\sqrt{r(1-r)}}\exp(1/(12j))
 <\frac{0.4}{\sqrt{v_*}\{1-1/(12\cdot40)\}}.
\]
Here $e^x\le(1-x)^{-1}$ for $0\le x<1$, proved by comparing their
power series. The endpoint gates ensure that the mode is interior.
This provides the concentration estimate at both $n-1$ and $n-2$.

The alternative concentration estimate obtained from the deleted
Berry--Esseen bounds retains \emph{both} hypotheses
$C_{n-1},C_{n-2}<0.4253$. Its scale factor $F_d$ must satisfy
\[
 F_d>0,\qquad F_d^2>\max\{(42/41)^3,42/40\}.
\]
The corresponding ratios decrease with $n$, so this endpoint gate is
sufficient for every $n\ge42$. A scale factor checked only at the older
floor need not satisfy it.

\subsection{The signed confinement step}

The signed deletion identities are algebraic identities at the actual
index $n-1$. Removing a pair leaves $n-2$ full summands. A centered
conditional noise $V$ satisfies $\E V_+=\E|V|/2$; replacing this signed
quantity by an unrelated absolute error would lose the factor needed
in the confinement estimate. Atoms and deterministic conditional laws
are covered by the same identities.

The exceptional pair event is strict. Its two candidate bounds have
the form
\[
 \Pp(\Delta>\tau)\le
 \min\left\{\frac{s}{s+p\tau^2},\frac{d_Y}{d_Y+\tau}\right\},
 \qquad \tau>0,
\]
with the conditional mean, variance and one-sided support assumptions
specified in the confinement derivation. The first is a one-sided
second-moment bound; the second uses the same inner gap. The interval
geometry additionally requires
\[
 1-w_{\rm strip}-2d_0>0,\qquad 1-2d_0-2\tau_+>0,
\]
as well as the disjointness and buffer gates in that derivation. These
conditions guarantee widths strictly below $1-d_Y$. The upper-ray
contradiction is a strict inequality of the form
$L-w_{\rm strip} U-\mathrm{error}>0$. The lower-ray argument also needs a
positive slope, so that a check at one endpoint excludes the whole
remaining ray inside the previously acquired strip. The new strip is
an output of this reasoning.

The step uses one strict phase resource,
\begin{equation}\label{ms:uniform-strict-resource}
 ne\rho+n\phi(0)D/6+n^{3/2}L_{\rm leak}<\mathcal B(H),
 \qquad e=C-\CE\ge0,\quad L_{\rm leak}\ge0.
\end{equation}
For a phase/noise cell, write its remaining resource as $R-nP$ with
$P\ge0$. Since $n\ge N$, it is at most $R-NP$. If the latter is
nonpositive, strict inequality~\eqref{ms:uniform-strict-resource}
excludes the cell. Otherwise a proved coefficient $\alpha>0$ gives
$nd_Y<(R-NP)/\alpha$. The error terms to which the index floor is
applied have the form
\[
 (zR+A)/n-zP+B\,\mathrm{bad},\qquad zR+A>0.
\]
The displayed positivity is what licenses $1/n\le1/N$; a formal
substitution without it would be invalid.

A sharper candidate uses the tangent inequality
\[
 \frac{x}{N\tau+x}\le
 \frac{N\tau}{(N\tau+t)^2}x+
 \frac{t^2}{(N\tau+t)^2},\qquad x,t\ge0,
\]
whose right side minus left side is
\[
 \frac{N\tau(x-t)^2}{(N\tau+x)(N\tau+t)^2}\ge0.
\]
The two error bounds are each valid with nonnegative coefficients and
the same resource simplex. Taking their minimum is then legitimate.

\subsection{Uniform assembly and verified tail closure}

The resulting extension has the following precise meaning. Retain the
attainment and contact hypotheses, the complete Fourier and signed
phase bounds, the ordered actual-law acquisition, all strict confinement
gates just described, the signed binomial normalization, and the full
positive-noise comparison. Replace all 65 finite support invocations
by Section~\ref{ms:support}. Retain the weighted invariant $nD\le42d_D$;
Section~\ref{ms:geometry-counts} then gives the joint plane at the same
floor. On each noise band, contract the support before using the new
amplitude in a subsequent estimate.

With every selected exact predicate verified at these inputs, the argument
proceeds at every actual $n\ge42$: signed entry, outer-panel exclusion,
central or left support contraction, signed confinement, joint-resource
export, noise-profile acquisition, and local count comparison. Closed
phase and noise seams and all admissible counts remain part of the
cover. The local signed cubic correction uses the actual upper score
when its coefficient is negative. The combined inverse-index correction
is clamped at zero before the floor substitution. At positive noise the comparison excludes the maximizing score; at zero
noise the finite-binomial comparison is used.

In the public verification release (\texttt{full\_pass: "ESTABLISHED"}), the
finite inequalities and analytic bounds are verified across the floor-40 role
(\texttt{current-C4042}), the 183-task finite collector role for $43\le n\le48$
(\texttt{current-C4348}), and the 55-predicate effective-tail suite for
$n\ge49$ (\texttt{independent-tail49}), including the ordered support
contractions, both deleted indices, the signed phase cells and all
positive-noise and zero-noise comparisons.
